\chapter{讨论与局限}

本文当前稿的第一个局限是 ground truth 边界。虽然 \capstone\ 为大规模指令扫描提供了现实可用的 ISA 视图，但它并不等价于 authoritative semantics。因此，凡是 decoder-related 的结论，都必须结合第二解码器、人工 triage 与后续实机验证理解。也正因如此，本文在 Layer A 结果解释中始终把当前输出视为 triage 证据，而不是直接上升为硬件事实。

第二个局限是分类输出尚未闭环。受控实验已经证明，现有实现能够稳定表达部分路径，但 illegal-at-test 与 known-disassembly exec-fault 这两类样例还没有稳定进入 artifact 日志。这个缺口既是当前实现的边界，也为后续工作指明了直接方向：在不破坏现有日志结构的前提下，补齐分类器对这些路径的落盘能力。

第三个局限是用户态方法的天然视野限制。用户态扫描非常适合做大范围差异审计，但它天然看不到全部处理器状态空间，尤其难以覆盖深层特权状态、缓存与 TLB 敏感路径，以及强状态依赖行为。因此，“未发现新隐藏指令”不应被解释为“处理器完全正确”；它最多说明，在当前子空间、当前环境和当前口径下，没有观察到足以支持该主张的证据。

第四个局限是 Layer B 尚未完成。当前论文已经可以围绕 Layer A 与 Layer C 形成一版结构完整的中文初稿，但只要没有真实硬件结果，任何关于实机可复现性和 Layer A/Layer B 差异归因的结论都必须保留为占位。这一点不仅影响实验章，也影响引言和结论中的论证强度。

最后，生成器的边界同样需要写清。本文并不声称配置生成器可以自动替代 ISA 工程，而是声称它能够把 bring-up 的高重复劳动前移到规格描述和基线生成阶段。对 \memcage\ 这样强依赖 ISA 语义的高性能后端来说，后续人工 hardening 仍然不可避免。
